% ── Abstract ─────────────────────────────────────────────────────────
% Target: 150–200 words. [R] numbers filled after real-footage ablation.
% ─────────────────────────────────────────────────────────────────────
Perception pipelines for urban traffic surveillance typically rely on
fixed detection thresholds and a single linear motion model, and both
assumptions break down in motorcycle-dominant mixed traffic: dense
weaving maneuvers fragment tracks, and static thresholds produce
persistent false-positive accident alerts.
%
We propose a context-aware adaptive perception pipeline with three
components: (C1)~a lightweight context estimator that reads traffic
density, time of day, image brightness, and recent false-positive rate,
and adjusts detector confidence, tracker buffer length, and motion-model
selection online with hysteresis-stabilized switching; (C2)~a
class-conditioned motion model that tracks motorcycles with a constant
turn-rate and velocity (CTRV) extended Kalman filter while other vehicle
classes retain the standard linear filter; and (C3)~a dual-task track
reliability score that serves high-precision accident alerting and
high-recall multi-agent signal-control observations from a single
pipeline via two task-specific thresholds.
%
On hand-annotated ground truth from a self-collected Vietnamese urban
dataset, C2 improves motorcycle identity preservation over a
YOLO11+ByteTrack baseline for $0.85$\,ms of additional per-frame tracker
latency, and improves both detection precision and recall simultaneously
--- something no per-frame confidence threshold can do.
We quantify the sensitivity of the identity gain to residual annotation
error rather than reporting a single figure.
We further show that C2's components are not separable: uncertainty-gated
re-association \emph{alone} degrades identity, and only becomes effective
once confidence-scaled measurement noise makes the filter's uncertainty
estimate trustworthy.
We report negatively on the remaining components --- context adaptation
does not improve on a calibrated fixed threshold in the single traffic
regime our annotated segment covers, and the reliability score is too
weakly discriminative (AUC $0.609$) to filter usefully once C2 has
removed most spurious tracks --- and identify what each would require to
be validated.
